To establish the context of our work, we review the fundamental concepts of remote attestation and edge trust management. Remote attestation allows a verifier to measure the software or firmware state of a target device (prover) and determine its integrity. Typically, this process utilizes a hardware-based Root of Trust, such as a Trusted Platform Module (TPM) or a hardware security module (HSM) embedded within the processor.

Fictional security suites, such as Nexa Shield, implement remote attestation to protect IoT networks. However, standard attestation is a point-to-point operation. If a single verifier becomes compromised, the trust status of the entire network is jeopardized. Decentralizing this process requires a distributed consensus mechanism. Prior attempts, such as Vertex-Trust~\cite{vertextrust2022}, utilize a permissioned blockchain. While Vertex-Trust provides decentralized trust, it requires all IoT nodes to participate in ledger consensus, which imposes significant energy and memory overhead. Synapse addresses these drawbacks by segregating the trust measurement phase from the ledger anchoring phase, utilizing edge-localized consensus groups.
